\section{Declaración de contribución individual}
\label{sec:contribuciones}

El trabajo fue desarrollado de manera colaborativa por los tres
integrantes del Grupo 5. Para organizar las responsabilidades, el
análisis se dividió en tres bloques temáticos coherentes: preparación y
exploración de los datos; inferencia frecuentista y regresión lineal; y
modelamiento predictivo e inferencia bayesiana.

La distribución porcentual representa la responsabilidad principal
asumida por cada integrante sobre el desarrollo, validación y
documentación de su bloque. Esta asignación no excluye la revisión
cruzada ni la participación conjunta en la interpretación de los
resultados.

La contribución total declarada es:

\[
40\,\% + 30\,\% + 30\,\% = 100\,\%.
\]

La Tabla~\ref{tab:contribuciones-individuales} resume la distribución
general de responsabilidades.

\begin{table}[H]
    \centering
    \caption{Distribución de las contribuciones individuales.}
    \label{tab:contribuciones-individuales}

    \small
    \begin{tabularx}{\textwidth}{
        @{}
        >{\raggedright\arraybackslash}p{3.1cm}
        >{\raggedright\arraybackslash}p{2.7cm}
        X
        @{}
    }
        \toprule
        Integrante
        & Bloque y participación
        & Responsabilidades principales \\
        \midrule

        Henry Sánchez Alvarado
        &
        \textbf{Bloque 1}

        Preparación, exploración e incertidumbre

        \textbf{40\,\%}
        &
        Organización general del flujo reproducible; documentación y
        carga del dataset; elaboración del diccionario de variables;
        validación, limpieza y recodificación de la respuesta; análisis
        exploratorio y estadística descriptiva; ajuste paramétrico del
        monto; desarrollo de las visualizaciones iniciales; estimación
        de intervalos de confianza mediante Wilson, bootstrap y
        distribución \(t\); integración de la estructura general del
        informe y revisión de la correspondencia entre notebook,
        resultados e informe.
        \\

        \addlinespace

        Armando Castro Chaupis
        &
        \textbf{Bloque 2}

        Inferencia frecuentista y regresión lineal

        \textbf{30\,\%}
        &
        Formulación y evaluación de las hipótesis frecuentistas;
        aplicación de la prueba chi-cuadrado y cálculo de la V de
        Cramér; comparación de la duración mediante la prueba de Welch y
        la \(d\) de Cohen; comparación del monto mediante
        Mann--Whitney y tamaños de efecto; especificación e
        interpretación del modelo OLS; evaluación de multicolinealidad,
        normalidad, homocedasticidad e influencia; comparación con
        errores robustos HC3 y validación cruzada del modelo lineal.
        \\

        \addlinespace

        Alex Segura Núñez
        &
        \textbf{Bloque 3}

        Predicción, costos e inferencia bayesiana

        \textbf{30\,\%}
        &
        Especificación y ajuste de la regresión logística; preparación
        de la partición estratificada de entrenamiento y prueba;
        interpretación de coeficientes y odds ratios; evaluación mediante
        ROC--AUC, KS, precisión, recall, especificidad, F1 y matrices de
        confusión; validación cruzada del clasificador; comparación de
        los umbrales \(0.50\) y \(1/6\) bajo la matriz de costos \(5:1\);
        desarrollo del modelo Beta--Binomial; análisis de sensibilidad
        al prior; implementación y diagnóstico de
        Metropolis--Hastings; integración de los resultados predictivos
        y bayesianos.
        \\

        \bottomrule
    \end{tabularx}
\end{table}

\subsection{Responsabilidades compartidas}

Aunque cada integrante asumió la responsabilidad principal de un bloque,
el trabajo no se desarrolló de manera aislada. Los tres miembros
participaron conjuntamente en:

\begin{itemize}
    \item definición y revisión de la pregunta de investigación;

    \item discusión de los supuestos estadísticos y del alcance del
    análisis;

    \item interpretación de los resultados desde la perspectiva del
    riesgo crediticio;

    \item revisión de la coherencia entre el notebook, las tablas CSV,
    las figuras y el informe;

    \item revisión crítica de las conclusiones, limitaciones y
    recomendaciones;

    \item preparación de la exposición y validación de la versión final
    del proyecto.
\end{itemize}

\subsection{Responsabilidad académica}

Cada integrante revisó el bloque bajo su responsabilidad y confirmó que
comprende:

\begin{itemize}
    \item la procedencia, estructura y codificación de los datos;

    \item las transformaciones y variables derivadas utilizadas;

    \item los métodos estadísticos y computacionales implementados;

    \item la interpretación de los coeficientes, intervalos y métricas;

    \item los supuestos y limitaciones de las pruebas y modelos;

    \item la diferencia entre asociación estadística, predicción y
    causalidad.
\end{itemize}

La responsabilidad académica sobre el contenido, los resultados, las
interpretaciones y las conclusiones corresponde exclusivamente a los
integrantes del Grupo 5. El uso de herramientas computacionales y de
inteligencia artificial se documenta en el
Anexo~\ref{anexo:uso-ia}.